% ---------------------------------------------------------------------------
% SECTION 4: ASYMPTOTIC THEORY (statements; proofs in Appendix A)
% ---------------------------------------------------------------------------
\section{Asymptotic Theory}\label{sec:theory}

This section states the theoretical properties of the A-SIMEX
estimator; complete proofs are collected in the Supplementary Material.
Throughout, $\|\cdot\|$ is the Euclidean norm,
$\rightarrow_p$ denotes convergence in probability,
$\rightarrow_d$ convergence in distribution, and generic positive
constants are written $C,c$.

\subsection{Setup and Assumptions}\label{sec:assumptions}

\begin{assumption}[Model and regularity]\label{ass:model}
\leavevmode
\begin{enumerate}\itemsep2pt
\item[(a)] $Y_1,X_1,\beps_1,\dots,Y_n,X_n,\beps_n$ are i.i.d.;
      $\E[Y_i\mid X_i]=\exp(X_i^\top\bet_0)$ with $\bet_0$ in the
      interior of the compact convex set $\Theta\subset\R^p$.
\item[(b)] $\E\|X_i\|^4<\infty$, $\E|Y_i|^2<\infty$, and
      $\E[X_iX_i^\top]$ is positive definite.
\item[(c)] The error satisfies Assumption~\ref{ass:cme}.
\end{enumerate}
\end{assumption}

\begin{assumption}[Pseudo-design non-collinearity]\label{ass:pseudo}
For every $\delta$ in a neighbourhood of $[-1,\lambda_{\max}+1]$ and
every $\bet$ in a neighbourhood of $\bet_0$,
\[
  \btH(\bet,\delta)
  := \E\Bigl[\exp\!\bigl(\bet^\top(X_i+\beps_\delta)\bigr)
      (X_i+\beps_\delta)(X_i+\beps_\delta)^\top\Bigr]
\]
is positive definite, where
$\beps_\delta\sim N(\bm 0,\delta\bSig_\epsilon)$; in particular
$\inf_{\bet,\delta}\lambda_{\min}\bigl(\btH(\bet,\delta)\bigr)\ge c_H>0$.
\end{assumption}

\begin{assumption}[Simulation and grid]\label{ass:simgrid}
The number of simulation replicates satisfies $B=B(n)\to\infty$. The
grid $\Lambda_n$ has $K=K(n)$ equally spaced points on
$[0,\lambda_{\max}]$ with $\lambda_{\max}\le\bar\lambda<\infty$ fixed,
$K\to\infty$, and $K^2\log n=o(n)$.
\end{assumption}

\begin{assumption}[Extrapolation degree; consistency
regime]\label{ass:degree}
The degree grid is $\{1,\dots,d_{\max}(n)\}$ with
$d_{\max}(n)\to\infty$, $d_{\max}(n)\le K(n)/4$, and the selected
degree $\hat d_n$ satisfies $\hat d_n\to\infty$ in probability and
$d_{\max}(n)^3/K(n)\to 0$.
\end{assumption}

\begin{assumption}[Selection stability; central limit
regime]\label{ass:cv}
The CV procedure selects $(\hat\lambda_{\max,n},\hat d_n)$ from a
fixed finite grid $\mathcal G$ with
$\Prob\bigl((\hat\lambda_{\max,n},\hat d_n)=(\lambda^*,d^*)\bigr)\to1$
for some $(\lambda^*,d^*)\in\mathcal G$.
\end{assumption}

\begin{remark}\label{rem:cvreconcile}
Assumptions~\ref{ass:degree} (degree grows: consistency) and
\ref{ass:cv} (degree eventually constant: the CLT) describe
complementary regimes; Theorem~\ref{thm:grow} reconciles them under an
explicit window condition on a $c\log n$ degree growth rate.
\end{remark}

\begin{remark}\label{rem:assumptions}
Three of these conditions correct or replace claims made in informal
treatments. First, no condition of the form $B/\sqrt n\to\infty$ (with
$B/n\to0$) is needed: the simulation component of the estimator is
$O_p(B^{-1/2})$ at the $\sqrt n$ scale, so $B\to\infty$ suffices for
all results. Second, a fixed polynomial degree cannot deliver
consistency in general: the extrapolation approximation error does
not vanish with $\lambda_{\max}$, which is why
Assumption~\ref{ass:degree} lets the CV-selected degree grow;
real-analyticity of the population curve
(Lemma~\ref{lem:curve}) then makes the approximation error decay
geometrically in the degree. Third, selection stability
(Assumption~\ref{ass:cv}) replaces ``convergence with positive
probability'' of the selected hyperparameters, which is too weak for a
central limit theorem; the cross-validation criterion supports both
regimes under verifiable conditions on the population risk
(Lemma~\ref{lem:cv}).
\end{remark}

\subsection{The Population SIMEX Curve}\label{sec:population}

Write $\delta=1+\lambda$ for the total error-variance multiplier and
$\beps_\delta=\beps_i+\sqrt\lambda\,\bu_i\sim N(\bm 0,
\delta\bSig_\epsilon)$ for the augmented error. Define the population
objective and pseudo-true parameter
\begin{equation}\label{eq:popobj}
  Q(\bet;\delta)
  =\E\Bigl[Y_i\,\bet^\top(X_i+\beps_\delta)
      -\exp\bigl(\bet^\top(X_i+\beps_\delta)\bigr)\Bigr],
  \qquad
  \beta(\delta)=\argmax_{\bet\in\Theta}Q(\bet;\delta),
\end{equation}
and write $\beta(\lambda):=\beta(1+\lambda)$ for the population SIMEX
curve.

\begin{lemma}[Population representation, analyticity, and the
extrapolation target]\label{lem:curve}
Grant Assumptions~\ref{ass:model} and~\ref{ass:pseudo}. Then:
\begin{enumerate}\itemsep2pt
\item[(i)] $Q$ depends on $\delta$ only through
      \[
        Q(\bet;\delta)\;=\;\E\bigl[Y_iX_i^\top\bet\bigr]
        \;-\;\E\Bigl[\exp\bigl(X_i^\top\bet\bigr)\Bigr]\,
        \exp\!\Bigl(\tfrac{\delta}{2}\,\bet^\top\bSig_\epsilon\bet\Bigr),
      \]
      hence $Q(\cdot\,;\delta)$ is real-analytic in
      $(\bet,\delta)$ on
      $\operatorname{int}\Theta\times(-1-\eta,\bar\lambda+1)$ for some
      $\eta>0$.
\item[(ii)] The maximizer $\beta(\delta)$ exists, is unique, lies in
      the interior of $\Theta$, and is real-analytic in $\delta$ on a
      neighbourhood of $[-1,\bar\lambda+1]$.
\item[(iii)] $\beta(-1)=\beta(0)=\bet_0$; moreover
      $\beta(\delta)=\bet_0+\delta\,\tilde\bet+O(\delta^2)$ as
      $\delta\downarrow0$, where, writing
      $\bm g(\bet)=\E\bigl[e^{X_i^\top\bet}X_i\bigr]$ and
      $M(\bet)=\E\bigl[e^{X_i^\top\bet}\bigr]$,
      \begin{equation}\label{eq:bias1}
        \tilde\bet
        =-\btH(\bet_0,0)^{-1}\Bigl[
          \tfrac12\bigl(\bet_0^\top\bSig_\epsilon\bet_0\bigr)\,
          \bm g(\bet_0)\;+\;\bSig_\epsilon\bet_0\,M(\bet_0)
        \Bigr],
      \end{equation}
      the leading-order bias of the naive estimator. For $p=1$,
      $X\sim N(0,1)$:
      $\tilde\beta=-\sigma^2_\epsilon\beta_0
      (1+\beta_0^2/2)/(1+\beta_0^2)<0$ for $\beta_0>0$.
\end{enumerate}
\end{lemma}

Two remarks on Lemma~\ref{lem:curve}. First, part~(iii) is the correct
form of the intuition that ``at $\lambda=-1$ the augmented covariate
equals $X$'': the simulated noise is independent of the observed error
and can never be subtracted off; instead, the augmented covariate at
level $\lambda$ is $X+\beps_\delta$ with \emph{total} error variance
$\delta\,\bSig_\epsilon$, and the pseudo-true parameter depends on
$\delta$ only through this total variance, so $\lambda=-1$
corresponds to a pseudo-population with no error at all, giving
$\beta(-1)=\bet_0$ exactly. Second, part~(i) is what drives every
approximation result below: the Gaussian tilt makes the population
curve real-analytic, so polynomial approximation error decays
geometrically in the degree (Lemma~\ref{lem:approx}).

\subsection{Uniform Convergence of the Empirical Curve}\label{sec:unif}

Let
$Q_n(\bet;\lambda,b)=n^{-1}\sum_{i=1}^n\bigl[
Y_i\,W_i^{(b)}(\lambda)^\top\bet
-\exp\{W_i^{(b)}(\lambda)^\top\bet\}\bigr]$ with
$W_i^{(b)}(\lambda)=W_i+\sqrt\lambda\,U_{ib}$ and
$\{U_{ib}\}$ i.i.d.\ $N(\bm 0,\bSig_\epsilon)$ across $i$ and $b$;
$\hat\beta_b(\lambda)=\argmax_{\bet\in\Theta}Q_n(\bet;\lambda,b)$;
$\bar\beta(\lambda)=B^{-1}\sum_b\hat\beta_b(\lambda)$; and
$\bm H(\lambda):=\btH(\beta(\lambda),1+\lambda)$.

\begin{lemma}[Uniform M-estimation over grid and
replicates]\label{lem:unif}
Grant Assumptions~\ref{ass:model}--\ref{ass:simgrid}. Then
\[
  \max_{0\le j<K}\ \max_{1\le b\le B}
  \sup_{\bet\in\Theta}\bigl|Q_n(\bet;\lambda_j,b)-Q(\bet;1+\lambda_j)\bigr|
  \;\xrightarrow{p}\;0,
\]
and consequently
$\max_{j}\|\bar\beta(\lambda_j)-\beta(\lambda_j)\|
=O_p(n^{-1/2})$ up to polylogarithms, with each
$\hat\beta_b(\lambda_j)$ well defined and unique with probability
tending to one.
\end{lemma}

\begin{lemma}[Extrapolation approximation error]\label{lem:approx}
Grant Assumptions~\ref{ass:model}--\ref{ass:simgrid}. There exist
constants $C,\rho>1$ depending only on the law of $(X_i,Y_i)$,
$\bSig_\epsilon$, $\bar\lambda$ and the grid such that for all $d\ge1$,
\[
  \inf_{\deg(q)\le d}\ \sup_{\lambda\in[-1,\bar\lambda]}
  \bigl\|q(\lambda)-\beta(\lambda)\bigr\|
  \;\le\;C\rho^{-d}.
\]
Consequently, under Assumption~\ref{ass:degree}, the best
degree-$\hat d_n$ approximation error is $o_p(n^{-1/2})$ whenever
$d_{\max}(n)\ge c\log n$ for large enough $c$.
\end{lemma}

\begin{lemma}[Cross-validation and degree
growth]\label{lem:cv}
Let $\widehat R_n(\lambda_{\max},d)$ be the CV criterion
\eqref{eq:cv} over the finite grid
$\mathcal G_n=\Lambda_n^{\max}\times\{1,\dots,d_{\max}(n)\}$ and
$R_n$ the population CV risk. If
(i)~$\sup_{g\in\mathcal G_n}|\widehat R_n(g)-R_n(g)|=o_p(1)$;
(ii)~$R_n(\lambda_{\max},d)\ge\rho^{-2d}(1-o(1))$ uniformly over
$d\le d_0$ for every fixed $d_0$; and (iii)~$R_n$ is minimized over
$d$ at some $\bar d_n\to\infty$, then the CV-selected degree
$\hat d_n\to\infty$ in probability. If moreover the grid contains a
fixed pair with a strict population minimum separated by more than
$2\sup_g|\widehat R_n(g)-R_n(g)|=o_p(1)$, then
Assumption~\ref{ass:cv} holds.
\end{lemma}

\subsection{Consistency}\label{sec:consistency}

\begin{theorem}[Consistency of \asimex{}]\label{thm:cons}
Grant Assumptions~\ref{ass:model}--\ref{ass:degree}. Then
$\hat\beta_{\asimex}\xrightarrow{p}\bet_0$. If
$d_{\max}(n)\ge c\log n$ eventually, then
$\|\hat\beta_{\asimex}-\bet_0\|=O_p(\rho^{-\hat d_n})
+O_p\bigl(\sqrt{d_{\max}(n)/K(n)}\,n^{-1/2}\bigr)+O_p(n^{-1/2})$.
\end{theorem}

\subsection{Asymptotic Normality and the Exact Variance
Decomposition}\label{sec:an}

For the central limit theorem we work in the fixed-hyperparameter
regime of Assumption~\ref{ass:cv}. To lighten notation write
$d=d^*$, $\lambda_{\max}=\lambda^*$, and
$\bphi(\lambda)=(1,\lambda,\dots,\lambda^d)^\top$. Define the
per-replicate influence score at grid point $\lambda_j$,
\[
  \bm m_{ib}(\lambda_j)
  =\Bigl[Y_i-\exp\bigl(\beta(\lambda_j)^\top W_i^{(b)}(\lambda_j)\bigr)\Bigr]
   W_i^{(b)}(\lambda_j),
\]
whose conditional mean given $(Y_i,X_i,\beps_i)$ does not depend on
$b$ and has zero unconditional mean at $\beta(\lambda_j)$; and the
matrices
\begin{align}
  \btV^{\mathrm{data}}
  &=\Cov\Bigl(\E\bigl[\bm m_{ib}(\lambda_j)\mid
        Y_i,X_i,\beps_i\bigr]\Bigr)_{j,k=0}^{K-1}
   \in\R^{Kp\times Kp},\label{eq:vdata}\\
  \btV^{\mathrm{sim}}
  &=\Cov\Bigl(\bm m_{ib}(\lambda_j)-\E[\bm m_{ib}(\lambda_j)\mid
        Y_i,X_i,\beps_i]\Bigr)_{j,k}
   \in\R^{Kp\times Kp},\label{eq:vsim}
\end{align}
the covariance contributed by the data and by the simulated errors,
stacked across the grid. Let
$\btH=\operatorname{blockdiag}(\bm H(\lambda_0),\dots,
\bm H(\lambda_{K-1}))$,
$\bm M=\bI_{p}\otimes[(\bPhi^\top\bPhi)^{-1}\bPhi^\top]$ the
least-squares extrapolation map, and
$\bm v_{-1}:=\bI_p\otimes\bphi(-1)$, so that
$\hat\beta_{\asimex}=\bm v_{-1}^\top\bm M\bar\beta_\Lambda$ with
$\bar\beta_\Lambda=(\bar\beta(\lambda_0)^\top,\dots,
\bar\beta(\lambda_{K-1})^\top)^\top$.

\begin{theorem}[Asymptotic normality of \asimex{}]\label{thm:an}
Grant Assumptions~\ref{ass:model}--\ref{ass:pseudo} and~\ref{ass:cv},
with $B\to\infty$ and fixed $K$, $\lambda_{\max}$, $d$. Then
\[
  \sqrt n\bigl(\hat\beta_{\asimex}-\bet_0\bigr)
  \;\xrightarrow{d}\;N\bigl(\bm 0,\bSig_{\asimex}\bigr),
\]
with the exact decomposition
\begin{equation}\label{eq:decomp}
  \bSig_{\asimex}
  \;=\;\bSig_{\mathrm{data}}\;+\;\tfrac1B\,\bSig_{\mathrm{sim}},
  \qquad
  \bSig_{\mathrm{data}}=\bm v_{-1}^\top\bm M\,
    \btH^{-1}\btV^{\mathrm{data}}\btH^{-1}\bm M^\top\bm v_{-1},
\end{equation}
and $\bSig_{\mathrm{sim}}$ defined analogously with
$\btV^{\mathrm{sim}}$ in place of $\btV^{\mathrm{data}}$. In
particular the simulation component vanishes for any rate
$B\to\infty$, including $B=o(n)$.
\end{theorem}

\begin{remark}[Correction to the additive
decomposition]\label{rem:decomp}
The asymptotic variance is the propagated influence covariance
\eqref{eq:decomp} of the pseudo-estimators at the grid points. An
additive split $\bSig_{\mathrm{param}}+\bSig_{\mathrm{extrap}}$ with
$\bSig_{\mathrm{param}}=\bm I_W(\bet_0)^{-1}$ (the information of the
naive working model) does not hold: the A-SIMEX estimator is a fixed
weighted average $\sum_jc_j\bar\beta(\lambda_j)$ of the grid-point
estimators with $\sum_jc_j=1$, the grid-point influence functions are
strongly (asymptotically perfectly, for small error) correlated, and
the naive information enters only through this correlation structure.
The correct additive split is between data and simulation sources of
randomness, with the simulation term exactly $O(B^{-1})$.
\end{remark}

\begin{corollary}[Variance estimation]\label{cor:var}
Under the conditions of Theorem~\ref{thm:an}, the plug-in estimator
described in Section~\ref{sec:var} (sample-analogue
$\bm H_j$ from the replicate-level information matrices, and empirical
covariance matrices of the estimated conditional influence means and
residuals) is consistent for $\bSig_{\asimex}$; equivalently, the
pairs-bootstrap (redrawing both data and simulated errors) is
consistent.
\end{corollary}

\subsection{The Growing-Degree Regime}\label{sec:growing}

Map $[0,\bar\lambda]$ affinely to $[-1,1]$ by
$t(\lambda)=2\lambda/\bar\lambda-1$, so that $\lambda=-1$ maps to
$t_0=-1-2/\bar\lambda<-1$, and define the extrapolation conditioning
constant
\[
  \kappa:=|t_0|+\sqrt{t_0^2-1}
  =1+\tfrac{2}{\bar\lambda}
   +\sqrt{\bigl(1+\tfrac{2}{\bar\lambda}\bigr)^2-1}\;>\;1,
\]
the Bernstein parameter of $t_0$: by the Chebyshev extremal property,
any polynomial $q$ of degree $k$ satisfies
$|q(t_0)|\le\kappa^{k}\sup_{[-1,1]}|q|$ \citep{rivlin1969}.

\begin{lemma}[Extrapolation operator bounds]\label{lem:extrap}
Let $K\ge c_0d^2$ equispaced points on $[0,\bar\lambda]$,
$\{\varphi_k\}_{k=0}^{d}$ the polynomials orthonormal with respect to
the discrete inner product $K^{-1}\sum_jf(\lambda_j)g(\lambda_j)$, and
$L_{K,d}(h)=\sum_jc_jh(\lambda_j)$ the ``least-squares degree-$d$ fit
evaluated at $\lambda=-1$'' functional. Then:
(i)~$\sum_{k=0}^{d}\varphi_k(t_0)^2\le C\,d\,\kappa^{2d}$;
(ii)~$\sum_j|c_j|\le C\,d^{1/2}K\kappa^{d}$ and
$L_{K,d}(q)=q(-1)$ for every polynomial of degree $\le d$;
(iii)~for every $h$ real-analytic on a neighbourhood of
$[-1,\bar\lambda]$ with analytic continuation bounded by $M$ on the
Bernstein ellipse $E_\rho$,
$\bigl|L_{K,d}(h)-h(-1)\bigr|\le C M d^{3/2}K\kappa^{d}\rho^{-d}$.
\end{lemma}

\begin{theorem}[Growing-degree central limit theorem]\label{thm:grow}
Grant Assumptions~\ref{ass:model}--\ref{ass:simgrid} with
$K(n)\ge c_0d_n^2$ equispaced points and $B\to\infty$. Suppose the
selected degree satisfies $\hat d_n/d_n\xrightarrow{p}1$ for a
deterministic sequence $d_n=\lceil c\log n\rceil$, and that the
analyticity margin $\rho>\kappa$ holds with
$c\in\bigl(\tfrac{1}{2\log\rho},\tfrac{1}{2\log\kappa}\bigr)$, an
interval that is nonempty exactly when $\rho>\kappa$. Then
\[
  \sqrt n\bigl(\hat\beta_{\asimex}-\bet_0\bigr)
  \;\xrightarrow{d}\;N\bigl(\bm 0,\bSig_*\bigr),
  \qquad
  \bSig_*=\bm I_X(\bet_0)^{-1}\,\widetilde{\bm Z}\,
           \bm I_X(\bet_0)^{-1},
\]
where $\widetilde{\bm Z}$ is the value at $(-1,-1)$ of the analytic
continuation of the influence covariance
$\bm\Gamma(\lambda,\lambda')
=\Cov(\bm\psi_1(\lambda),\bm\psi_1(\lambda'))$ of the data-influence
process $\bm\psi(\lambda)=\E[\bm m_{ib}(\lambda)\mid Y_i,X_i,\beps_i]$,
and $\bm I_X(\bet_0)=\btH(-1)$. The simulation component contributes
$O(B^{-1})$, and the extrapolation bias is $o(n^{-1/2})$, so the limit
is centered at $\bet_0$. In particular
$\bSig_*\to\bSig_{\mathrm{oracle}}$ as $\sigma^2\to0$.
\end{theorem}

The window condition has a transparent reading: $\rho$ measures how
far the analytic continuation of the curve extends (bias side),
$\kappa$ how violently degree-$d$ extrapolation outside the grid
amplifies noise (variance side); growing-degree inference is feasible
precisely when the former dominates the latter. If $\rho\le\kappa$,
the fixed-degree central limit theorem (Theorem~\ref{thm:an}) is the
appropriate inference tool.

\subsection{Efficiency Relative to the Oracle}\label{sec:are}

Let $\bSig_{\mathrm{oracle}}=\bm I_X(\bet_0)^{-1}$ and define the
generalized asymptotic relative efficiency
$\mathrm{ARE}=\det(\bSig_{\mathrm{oracle}})/
\det(\bSig_{\asimex})$. Three information matrices must be
distinguished: the oracle information $\bm I_X(\bet_0)$; the
information $\bm I(t)$ of the observed-data model
$(Y_i,\widetilde W_t)$, $\widetilde W_t=X_i+\bm e_t$ with
$\bm e_t\sim N(\bm 0,t\bSig_0)$; and the information of the naive
\emph{working} model $\E[e^{W^\top\bet_0}WW^\top]$, a curvature of a
misspecified likelihood that plays no role in the correct theory.

\begin{lemma}[Information monotonicity under additive
noise]\label{lem:info}
Grant Assumption~\ref{ass:model} with $\bSig_0$ fixed. Then:
(i)~$\bm I(t_2)\preceq\bm I(t_1)$ for $0\le t_1\le t_2$; in
particular $\bm I(t)\preceq\bm I(0)=\bm I_X(\bet_0)$; (ii)~the
inequality is strict for $t>0$ provided $X_i$ is not a.s.\ constant.
Consequently
$\bSig_{\mathrm{oracle}}\preceq\bm I(t)^{-1}$ for all $t\ge0$, and by
the H\'ajek convolution theorem $\bm I(t)^{-1}$ is the best asymptotic
covariance achievable by any regular estimator based on
$(Y_i,\widetilde W_t)$ alone.
\end{lemma}

\begin{theorem}[Efficiency relative to the oracle]\label{thm:are}
Grant the conditions of Theorem~\ref{thm:an} with $B\to\infty$, and
suppose $\bSig_\epsilon=\sigma^2\bSig_0$ with $\bSig_0$ fixed and
$\sigma^2\downarrow0$, all else fixed. Write
$\bSig_{\asimex}(\sigma^2)$ for the asymptotic covariance of
Theorem~\ref{thm:an} at error scale $\sigma^2$. Then:
\begin{enumerate}\itemsep2pt
\item[(i)] \textbf{(First-order efficiency.)}
      $\bSig_{\asimex}(\sigma^2)\to\bSig_{\mathrm{oracle}}$ as
      $\sigma^2\to0$, for every fixed grid and degree: the limit is
      exact.
\item[(ii)] \textbf{(Explicit second-order term.)}
      $\bSig_{\asimex}(\sigma^2)
      =\bSig_{\mathrm{oracle}}+\sigma^2\,\dot{\bm C}_\beta+O(\sigma^4)$,
      where $\dot{\bm C}_\beta$ is the symmetric, explicitly
      computable matrix
      \begin{equation}\label{eq:cdot}
        \dot{\bm C}_\beta
        =\bm v_{-1}^\top\bm M\,\bm H_0^{-1}
        \Bigl(\dot{\btV}-\dot{\btH}\bm H_0^{-1}\bm V_0
              -\bm V_0\bm H_0^{-1}\dot{\btH}\Bigr)
        \bm H_0^{-1}\bm M^\top\bm v_{-1},
      \end{equation}
      with $\bm H_0=\bI_K\otimes\bm I_X(\bet_0)$,
      $\bm V_0=(\bm 1_K\bm 1_K^\top)\otimes\bm I_X(\bet_0)$, and
      \begin{align}
        \dot{\btH}_j
        &=\tfrac{\delta_j}{2}\bigl(\bet_0^\top\bSig_0\bet_0\bigr)\bm I_X
          +\delta_j\,\E\bigl[e^{\bet_0^\top X}
             (\tilde\bet_{\bSig_0}^\top X)XX^\top\bigr]
          +\bSig_0\bet_0\,\bm g_0^\top+\bm g_0\,\bet_0^\top\bSig_0,
          \label{eq:hdot}\\
        \dot{\btV}_{jk}
        &=\Cov\bigl(\dot{\bm\psi}_j,\bm\psi^{(0)}\bigr)
         +\Cov\bigl(\bm\psi^{(0)},\dot{\bm\psi}_k\bigr),
        \nonumber\\
        \dot{\bm\psi}_j
        &=-e^{\bet_0^\top X}
         \Bigl[\bigl(\delta_j\tilde\bet_{\bSig_0}^\top X
                +\tfrac{\lambda_j}{2}\bet_0^\top\bSig_0\bet_0\bigr)X
               +\lambda_j\bSig_0\bet_0\Bigr],
        \label{eq:vdot}
      \end{align}
      where $\delta_j=1+\lambda_j$,
      $\bm g_0=\E[e^{\bet_0^\top X}X]$, and
      $\bm\psi^{(0)}=YX-e^{\bet_0^\top X}X$ is the oracle influence
      function. In general $\dot{\bm C}_\beta$ is
      \emph{indefinite}: neither
      $\dot{\bm C}_\beta\succeq\bm 0$ nor
      $\dot{\bm C}_\beta\preceq\bm 0$ holds universally.
\item[(iii)] \textbf{(Determinant expansion.)}
      $\mathrm{ARE}=1+\sigma^2\kappa_\beta+O(\sigma^4)$ with
      $\kappa_\beta=-\operatorname{tr}
      (\bSig_{\mathrm{oracle}}^{-1}\dot{\bm C}_\beta)$ of no universal
      sign. The universal statements are: (a)~$\mathrm{ARE}\to1$ as
      $\sigma^2\to0$; (b)~by Lemma~\ref{lem:info}, no regular
      estimator based on the noisy covariates alone can improve on the
      oracle to first order; and (c)~\asimex{} attains this frontier
      to first order, its excess over $\bSig_{\mathrm{oracle}}$ being
      second-order with a bias--variance trade-off reflected in the
      sign of $\kappa_\beta$.
\end{enumerate}
\end{theorem}

\begin{remark}[Why a positive-semidefiniteness claim would be
false]\label{rem:psd}
Measurement error attenuates the signal but can also de-noise the
estimating directions: in the homoskedastic linear model
$Y=X\beta_0+e$, $W=X+\epsilon$, the naive slope's asymptotic variance
$\sigma_e^2/\E[W^2]$ is \emph{smaller} than the oracle's
$\sigma_e^2/\E[X^2]$, and the second-order variance term inherits this
negative sign; since A-SIMEX is a biased estimator, it is not subject
to the Cram\'er--Rao bound of the observed-data model. The universal
content is part~(iii): the oracle marks the first-order frontier,
A-SIMEX reaches it, and the second-order gap splits into the
deterministic extrapolation bias
$\sigma^2(\sum_jc_j\delta_j)\tilde\bet_{\bSig_0}$ (zero for suitable
grids) plus the symmetric variance correction \eqref{eq:cdot} of
indefinite sign.
\end{remark}

\begin{corollary}[Explicit scalar constants]\label{cor:scalar}
For $p=1$ the first-order naive bias is
$\tilde\beta=-\sigma^2\bigl[\tfrac12\beta_0^2\mu_1(\beta_0)
+\beta_0\mu_0(\beta_0)\bigr]/\mu_2(\beta_0)$ with
$\mu_r(\beta)=\E[e^{\beta X}X^r]$ (for $X\sim N(0,1)$,
$-\sigma^2\beta_0(1+\beta_0^2/2)/(1+\beta_0^2)$), and the
second-order ARE constant $\kappa_\beta$ is an explicit functional of
$\{\mu_r(\beta_0)\}_{r\le4}$ and the grid weights $\{c_j\}$, given in
the Supplementary Material (Section~S1).
\end{corollary}
